\section{ماژول swift\_svd\_act (Swift SVD with Activation Subspace Projection)}
\label{sec:module_swift_svd_act}

\subsection{معرفی و انگیزه}
این ماژول در چارچوب پژوهش حاضر پیاده‌سازی و ارزیابی شده و فاقد مقاله مستقل جداگانه در ادبیات تحقیق است.

تجزیه مقادیر منفرد ایستا (\lr{Static SVD}) به تنهایی قادر به تفکیک زیرفضاهای حامل اطلاعات حیاتی از مولفه‌های کم‌اهمیت در جریان پردازش پویای شبکه نیست. در مقابل، ثبت کامل تاریخچه فعال‌سازی‌ها در حافظه رم نیز به دلیل انفجار حجم تانسورها در مدل‌های عمیق امکان‌پذیر نمی‌باشد. برای دستیابی به حداکثر دقت بازسازی با کمترین سربار پردازشی، نیاز به الگوریتمی است که ماتریس کوواریانس فعال‌سازی‌های ورودی و خروجی را به صورت برخط تجمیع کرده و تصویر وزن‌ها را بر روی زیرفضای اصلی فعال‌سازی هدایت نماید.

ماژول \lr{swift\_svd\_act} با نام فارسی «فشرده‌ساز آگاه از فعال‌سازی Swift بر مبنای تصویر زیرفضا» نسخه پیشرفته و داده‌آگاه از فشرده‌ساز \lr{swift\_svd} است. این ماژول با استفاده از هوک‌های برخط (\lr{Forward Hooks})، ماتریس کوواریانس متقاطع ورودی-خروجی را در محاسبات با دقت مضاعف (\lr{float64}) انباشته کرده، پایه متعامد زیرفضای برانگیختگی را استخراج نموده و ماتریس وزن لایه را بر این زیرفضا تصویر می‌کند تا جهت‌های پرانرژی سیگنال با بالاترین دقت ممکن حفظ شوند.

\subsection{مبانی نظری و فرمول‌بندی ریاضی}
فرض کنید برای یک لایه خطی با وزن $\mathbf{W} \in \mathbb{R}^{d_{\text{out}} \times d_{\text{in}}}$، جفت تانسورهای ورودی و خروجی روی $N$ توکن کالیبراسیون به ترتیب $\mathbf{X} \in \mathbb{R}^{N \times d_{\text{in}}}$ و $\mathbf{Y} \in \mathbb{R}^{N \times d_{\text{out}}}$ باشند.

\textbf{۱. انباشت برخط کوواریانس متقاطع (\lr{Online Cross-Covariance}):} \\
ماتریس کوواریانس متقاطع تجربی به صورت برخط و بدون ذخیره نمونه‌های منفرد در حافظه محاسبه می‌شود:
\begin{equation}
\mathbf{C}_{XY} = \frac{1}{N} \mathbf{Y}^\top \mathbf{X} \in \mathbb{R}^{d_{\text{out}} \times d_{\text{in}}}
\end{equation}
انباشت این تانسور در دقت مضاعف (\lr{torch.float64}) صورت می‌گیرد تا از خطاهای ناشی از سرریز و پدیده کم‌شدن دقت جلوگیری شود.

\textbf{۲. استخراج زیرفضای فعال‌سازی (\lr{Subspace Extraction}):} \\
تجزیه مقادیر منفرد بر روی ماتریس کوواریانس متقاطع انجام می‌شود:
\begin{equation}
\mathbf{C}_{XY} = \mathbf{U}_c \mathbf{\Sigma}_c \mathbf{V}_c^\top
\end{equation}
با برش به رتبه $r_{\text{act}} \le r$， پایه‌های قطعی فعال‌سازی خروجی و ورودی استخراج می‌گردند:
\begin{equation}
\mathbf{U}_{\text{act}} = \mathbf{U}_{c}[:, :r_{\text{act}}] \in \mathbb{R}^{d_{\text{out}} \times r_{\text{act}}}, \quad \mathbf{V}_{\text{act}} = \mathbf{V}_{c}[:, :r_{\text{act}}] \in \mathbb{R}^{d_{\text{in}} \times r_{\text{act}}}
\end{equation}

\textbf{۳. تصویر وزن بر زیرفضای فعال‌سازی و سنتز فاکتورها (\lr{Projection and Factor Synthesis}):} \\
ماتریس وزن اصلی بر زیرفضای دوجانبه فعال‌سازی تصویر می‌شود:
\begin{equation}
\mathbf{M} = \mathbf{U}_{\text{act}}^\top \mathbf{W} \mathbf{V}_{\text{act}} \in \mathbb{R}^{r_{\text{act}} \times r_{\text{act}}}
\end{equation}
سپس تجزیه نهایی مقادیر منفرد بر روی این ماتریس کوچک انجام می‌پذیرد:
\begin{equation}
\mathbf{M} = \mathbf{U}_m \mathbf{\Sigma}_m \mathbf{V}_m^\top
\end{equation}
فاکتورهای نهایی بازسازی وزن با فرمول‌بندی زیر سنتز می‌شوند:
\begin{align}
\mathbf{A} &= (\mathbf{U}_{\text{act}} \mathbf{U}_m[:, :r]) \operatorname{diag}(\sqrt{\mathbf{\sigma}_{m, 1:r}}) \in \mathbb{R}^{d_{\text{out}} \times r} \\
\mathbf{B} &= \operatorname{diag}(\sqrt{\mathbf{\sigma}_{m, 1:r}}) (\mathbf{V}_m^\top[:r, :] \mathbf{V}_{\text{act}}^\top) \in \mathbb{R}^{r \times d_{\text{in}}}
\end{align}
به طوری که $\mathbf{W} \approx \mathbf{A} \mathbf{B}$ بوده و خطای بازسازی به صورت تابع هزینه فروبنیوس $\mathcal{E}_{\text{rel}} = \frac{\|\mathbf{W} - \mathbf{A}\mathbf{B}\|_F}{\|\mathbf{W}\|_F + 10^{-8}}$ پایش می‌گردد. در صورت عدم همگرایی مناسب، الگوریتم به صورت خودکار به \lr{SVD} مستقیم سوئیچ می‌کند.

\subsection{معماری و پیاده‌سازی}
پیاده‌سازی این الگوریتم در فایل \lr{neuralops/compression\_engine/swift\_svd\_act.py} مستقر بوده و تحت شناسه‌های \lr{swift\_svd\_act}، \lr{swift\_svd} و \lr{swift\_svd\_compressor} در رجیستری سیستم ثبت شده است. نقاط قوت معماری آن عبارتند از:
\begin{itemize}
    \item \textbf{انباشت آنلاین بدون انفجار رم (\lr{SW-FIX-04}):} بهره‌گیری از هوک‌های موقت پیشرو که تانسورهای کوواریانس را مستقیماً محاسبه کرده و تانسورهای فعال‌سازی حجیم ورودی و خروجی را بلافاصله آزاد می‌کنند.
    \item \textbf{پاک‌سازی قطعی هوک‌ها در بلوک‌های \lr{try/finally} (\lr{SW-FIX-06}):} تضمین حذف هوک‌های ثبتی حتی در صورت بروز استثنا در جریان گذر رو به جلو کالیبراسیون.
    \item \textbf{پایداری عددی با \lr{float64} (\lr{SW-FIX-07}):} انجام محاسبات ضرب تنسوری کوواریانس در فضای اعداد اعشاری ۶۴ بیتی برای حفظ تقارن عددی.
    \item \textbf{استراتژی دومرحله‌ای بازگشت امن (\lr{Two-Tier Fallback}):} در صورتی که خطای ناشی از تصویر زیرفضا بیشتر از حد مجاز باشد، ابتدا تجزیه SVD مستقیم آزموده می‌شود؛ در صورت تداوم خطای بالا، وزن اصلی لایه بدون تغییر حفظ می‌گردد.
\end{itemize}

\subsection{نتایج تجربی}
ارزیابی‌های تجربی فشرده‌ساز آگاه از فعال‌سازی \lr{swift\_svd\_act} بر اساس اطلاعات پایگاه داده اعتبارسنجی یکپارچه \lr{thesis\_final/master\_verified\_results.json} در جدول \ref{tab:benchmark_swift_svd_act} منعکس گردیده است:

\begin{table}[H]
\centering
\small
\caption{نتایج بنچ‌مارک تجربی فشرده‌ساز آگاه از فعال‌سازی \lr{swift\_svd\_act} روی مدل‌های پایه}
\label{tab:benchmark_swift_svd_act}
\begin{tabular}{lccccc}
\toprule
\textbf{مدل پایه \lr{(Base Model)}} & \textbf{دقت \lr{(Accuracy)}} & \textbf{شاخص \lr{F1}} & \textbf{حافظه \lr{(Peak VRAM)}} & \textbf{توان \lr{(Throughput)}} & \textbf{وضعیت} \\
\midrule
\lr{DistilBERT} & \lr{90.48\%} & \lr{0.9077} & \lr{311.19 MB} & \lr{146,011.97 tok/s} & \lr{PASS} \\
\lr{RoBERTa} & \lr{72.02\%} & \lr{0.7724} & \lr{301.72 MB} & \lr{91,839.43 tok/s} & \lr{PASS} \\
\lr{GPT-2} & \lr{60.09\%} & \lr{0.6463} & \lr{637.29 MB} & \lr{66,501.09 tok/s} & \lr{PASS} \\
\lr{Qwen} & \lr{49.08\%} & -- & -- & -- & \lr{PASS} \\
\bottomrule
\end{tabular}
\end{table}

مقایسه نتایج جدول \ref{tab:benchmark_swift_svd_act} با نگارش بدون کالیبراسیون (\lr{swift\_svd}) تفاوت معنادار بهره‌گیری از تصویر زیرفضای فعال‌سازی را اثبات می‌کند: در مدل \lr{RoBERTa} دقت از \lr{50.92\%} به \lr{72.02\%} (جهش ۲۱.۱ درصدی) ارتقا یافته و وضعیت از \lr{PASS\_LOW} به قبولی قطعی \lr{PASS} تغییر یافته است. در مدل خودبازگشت \lr{GPT-2} نیز دقت با ارتقای ۹.۱۷ درصدی به \lr{60.09\%} رسیده و قبولی کامل \lr{PASS} احراز شده است. هم‌زمان، در مدل \lr{DistilBERT} دقت بسیار بالای \lr{90.48\%} با توان عملیاتی خیره‌کننده بیش از ۱۴۶ هزار توکن بر ثانیه و کمترین مصرف حافظه گرافیکی ($301.72$ مگابایت در روبرتا) تثبیت شده است. تمام چهار مدل در این ارزیابی وضعیت قبولی کامل \lr{PASS} را به دست آورده‌اند.

\subsection{محدودیت‌ها و تحلیل شکست}
تحلیل رفتار الگوریتم و شرایط مرزی نشان‌دهنده ابعاد زیر است:
\begin{itemize}
    \item \textbf{وابستگی کیفیت به نماینده بودن داده‌های کالیبراسیون:} دقت استخراج زیرفضا تا حد زیادی به تطابق محتوای متون کالیبراسیون با توزیع داده‌های آزمون بستگی دارد؛ در صورت انحراف دامنه داده‌ها (\lr{Domain Shift})، جهت‌های زیرفضا دچار انحراف گردیده و لایه ناچار به عقب‌نشینی به SVD مستقیم می‌شود.
    \item \textbf{سربار حافظه‌ای هوک‌ها در مدل‌های بسیار بزرگ:} در لایه‌های با طول توالی بسیار طولانی ($L > 4096$)، ضرب ماتریسی فعال‌سازی‌ها برای ساخت کوواریانس می‌تواند بار لحظه‌ای بر حافظه تحمیل کند که البته با سقف‌گذاری توکن‌ها در پیاده‌سازی کنترل شده است.
\end{itemize}
